% ECONOMICS_PROBLEM: {"id": "D71-9", "title": "The optimal randomized ordinal metric-distortion constant", "jel_code": "D71", "source_id": "OP-0227"}
% !TeX program = xelatex
\documentclass[11pt,a4paper]{article}
\usepackage[margin=23mm]{geometry}
\usepackage{fontspec}
\setmainfont{Latin Modern Roman}
\usepackage{amsmath,amssymb,mathtools}
\usepackage{xcolor,enumitem,booktabs,longtable,array,xurl,hyperref}
\definecolor{muted}{HTML}{53636C}
\hypersetup{colorlinks=true,urlcolor=blue,linkcolor=blue}
\setlength{\parindent}{0pt}
\setlength{\parskip}{5pt}
\newcommand{\R}{\mathbb R}
\newcommand{\E}{\mathbb E}
\newcommand{\N}{\mathbb N}
\newcommand{\ind}{\mathbf 1}
\newcommand{\Var}{\operatorname{Var}}
\newcommand{\Cov}{\operatorname{Cov}}
\newcommand{\supp}{\operatorname{supp}}
\DeclareMathOperator*{\argmin}{arg\,min}
\DeclareMathOperator*{\argmax}{arg\,max}
\newcommand{\entrymeta}[3]{{\sffamily\small\color{muted}\textbf{\texttt{#1}}\quad|\quad #2\par #3\par}}
\newcommand{\Problemref}[1]{\texttt{#1}}
\newcommand{\JELref}[2]{\texttt{#2} (JEL \texttt{#1})}
\begin{document}
\section*{The optimal randomized ordinal metric-distortion constant}
\textbf{Problem ID:} \texttt{D71-9}\quad \textbf{JEL:} \texttt{D71}\par
% BEGIN ECONOMICS STATEMENT
\label{problem:OP-0227}
% GAME_THEORY_PROBLEM: {"local_id": "GT-097", "original_number": 97, "kind": "P", "entry_kind": "statement_only", "published": false, "review_required": true, "category": "game_theory"}
\label{game:GT-097}
{\sffamily\small\color{muted}
\noindent\textbf{ID:} GT-097. \textbf{Category:} Randomized metric social choice.
\textbf{Status:} P; closing the remaining distortion gap is requested in the inspected source.
\par}
\subsection*{Definitions and mathematical statement}
For each $n\geq1,m\geq2$, let $V,C$ be labelled disjoint sets of those sizes and let $\mathcal P_{n,m}$ be the set of profiles of strict candidate rankings. A randomized ordinal rule is a family
\[
 R=(R_{n,m})_{n,m},\qquad R_{n,m}:\mathcal P_{n,m}\to\Delta(C),
 \quad\Delta(C)=\{q\in\mathbb R_{\geq0}^{C}:\sum_cq_c=1\}.
\]
It receives only the rankings. For a profile $P$, let $\mathcal D(P)$ be all pseudometrics on $V\sqcup C$ such that $a\succ_v b$ implies $d(v,a)\leq d(v,b)$, and put
$\operatorname{SC}_d(c)=\sum_vd(v,c)$.
Define
\[
 \begin{split}
 D(R)=\inf\{D\geq0:\ &\forall n,m\ \forall P\in\mathcal P_{n,m}\
                         \forall d\in\mathcal D(P),\\
 &\sum_{c\in C}R_{n,m}(P)(c)\operatorname{SC}_d(c)
      \leq D\min_{a\in C}\operatorname{SC}_d(a)\},
 \qquad D^*=\inf_R D(R).
 \end{split}
\]
Take $D(R)=+\infty$ if no finite constant works. Determine $D^*$, ideally by a matching uniform rule family and lower bound. The equivalent ratio supremum is restricted to positive optimal cost, together with the requirement that expected cost vanish whenever optimal cost is zero.

\subsection*{Short English statement}
What is the smallest worst-case expected distortion achievable from full ordinal rankings, uniformly over all election sizes?

\subsection*{Variants and scope boundaries}
The supremum includes ranking profiles and electorate sizes, not just metrics for one fixed election. This target imposes neither strategyproofness nor computational-efficiency restrictions. Restricting the candidate number, using partial elicitation, or requiring determinism changes the optimum. The metric is fixed before the rule's lottery is realized.

\subsection*{Source relationship and status}
The inspected source gives an upper guarantee of $2.13713$ and reports the lower bound $2.11264$, leaving a gap. Thus its stated bounds place this model's optimum between those reported constants; the catalog's $2.5$ progress figure is outdated for this source. The exact infimum and possible attainment are separate questions.

\subsection*{References}
\begingroup\footnotesize\raggedright
\noindent [S1] Ziyi Cai, Moses Charikar, Jabari Hastings, Prasanna Ramakrishnan, Kangning Wang, and Qilin Ye, \emph{Stable Voting Rules on the Edge of Optimal Metric Distortion} (2026), arXiv:2609.08259v1. Inspected locators: Section 2 (metric-election model and distortion), Theorem 1 and introductory bounds, Section 7 (remaining gap).
\url{https://arxiv.org/html/2609.08259v1}
\par\endgroup

\subsection*{Common conventions: Source mathematical conventions}

\textbf{Finite games.} For a finite set $A$, $\Delta(A)$ is its probability
simplex. Players choose independent mixed strategies in Nash equilibrium;
correlation is allowed only when explicitly part of the solution concept.
In a game with payoff functions $u_i$ and strategy profile $x$, an additive
$\varepsilon$ Nash equilibrium satisfies
\[
 u_i(x)\geq u_i(x_i',x_{-i})-\varepsilon
 \quad\text{for every player $i$ and every admissible deviation $x_i'$}.
\]
For numerical approximation, payoffs are normalized as locally specified.
An $\varepsilon$ payoff residual is distinct from distance at most
$\varepsilon$ to the set of exact equilibria. Point convergence, convergence
of empirical play, and convergence in distance to an equilibrium set are
also distinct.

\textbf{Computational representations.} Unless stated otherwise, a complexity
question takes rational data with binary encoding; polynomial time is measured
in total bit length. A fixed-accuracy polynomial algorithm is not necessarily
polynomial in $1/\varepsilon$, and neither is necessarily polynomial in
$\log(1/\varepsilon)$. Succinct games and value, demand, or sampling oracles
must declare their interfaces. Exact real solutions require an explicit
representation; an unspecified list of arbitrary real numbers is not a
finite computational output. A claimed algorithmic barrier is meaningful
only for its specified model and reductions.

\textbf{Dynamic games.} Histories, information, observation, and allowable
randomization are part of the model. A stationary strategy depends only on
the current state; a positional strategy in a deterministic graph game
chooses one action at each controlled vertex. Nash equilibrium at an initial
state and equilibrium after every history are different requirements.
An expectation of a pathwise $\liminf$ payoff, a $\liminf$ of expected averages,
a limiting expected average, and a uniform finite-horizon equilibrium are
different criteria. None is silently substituted for another.

\textbf{Allocation and mechanisms.} A complete allocation assigns every item
exactly once unless otherwise stated. Zero-valued goods, negative values,
capacity constraints, feasibility, lotteries, and copies of goods affect
fairness definitions and are specified locally. Dominant-strategy incentive
compatibility, Bayesian incentive compatibility, universal truthfulness,
and truthfulness in expectation impose different inequalities. Whenever
an objective ratio has zero denominator, its convention is stated or those
instances are excluded.

\textbf{Infinite-dimensional models.} Expectations, integrals, stochastic
equations, and optimization objectives require their local regularity,
integrability, filtrations, and admissible domains. A backward--forward
mean-field system is a boundary-value problem; it is not an initial-value
semiflow without an additional construction. Long-time, large-population,
and vanishing-noise limits require an order of limits and a convergence
topology. If a minimum need not be attained, the objective is its infimum
or supremum and the approximation requirement is stated separately.
% END ECONOMICS STATEMENT
\end{document}
